% chapters/state_of_art/technologies.tex
% Sección 2.2: Tecnologías y herramientas (~2-3 páginas)
% ============================================================================
%
% GUÍA: Describir las tecnologías utilizadas en el proyecto, justificando
%       la elección de cada una frente a alternativas. Para cada tecnología:
%       - Qué es y para qué sirve
%       - Por qué se eligió (ventajas frente a alternativas)
%       - Citar la documentación oficial o artículos relevantes
%
%   Organizar por subsecciones (sin numerar, usar \paragraph{}):
%
%   \paragraph{Python y Django}
%   - Python como lenguaje de propósito general, madurez, ecosistema.
%   - Django: framework web full-stack, ORM, admin, migraciones.
%   - Alternativas descartadas: Flask (microframework, menos estructura),
%     FastAPI (más moderno pero menos maduro), Spring Boot (Java).
%   - Citar: ~\cite{django2024,python2024}
%
%   \paragraph{Django REST Framework}
%   - Serialización, vistas genéricas, autenticación, permisos, paginación.
%   - Alternativas: Django Ninja, Tastypie.
%   - Citar: ~\cite{drf2024}
%
%   \paragraph{PostgreSQL}
%   - RDBMS robusto, soporte JSON, extensiones, rendimiento.
%   - Alternativas: MySQL/MariaDB, SQLite (no concurrente).
%   - Citar: ~\cite{postgresql2024}
%
%   \paragraph{Celery y Redis}
%   - Celery: cola de tareas distribuida para Python.
%   - Redis: broker de mensajes, caché.
%   - Alternativas: RabbitMQ (broker), Dramatiq (cola), Huey.
%   - Citar: ~\cite{celery2024,redis2024}
%
%   \paragraph{MinIO}
%   - Almacenamiento de objetos compatible con S3.
%   - Para PDFs, anotaciones de audio, imágenes rasterizadas.
%   - Alternativas: almacenamiento local, Amazon S3 directamente.
%   - Citar: ~\cite{minio2024}
%
%   \paragraph{Reconocimiento óptico de caracteres (OCR)}
%   TODO: Especificar qué motor de OCR usas (Tesseract, EasyOCR, PaddleOCR,
%         Google Vision, etc.). Describir brevemente cómo funciona y por qué
%         se eligió. Mencionar la arquitectura pluggable que permite cambiar
%         el motor sin afectar al resto del sistema.
%   - Citar: ~\cite{tesseract2024} o el motor elegido.
%
%   \paragraph{Docker y Docker Compose}
%   - Contenerización de servicios, reproducibilidad, despliegue.
%   - Citar: ~\cite{docker2024}
%
%   \paragraph{Otras herramientas}
%   TODO: Añadir según se vayan usando. Posibles candidatas:
%   - drf-spectacular (generación OpenAPI/Swagger)
%   - pytest / pytest-django (pruebas)
%   - Pillow (procesamiento de imágenes)
%   - WeasyPrint o reportlab (generación de PDFs rasterizados)
%   - Herramientas de CI/CD si se configuran
%   - Git, GitHub/GitLab
% ============================================================================

% TODO: Redactar la revisión de tecnologías.
